% English translation of checked_pages/esq_119_008.tex; source: cote 119, site 8, PDF page 9, typescript page 6.
% AI-translated by GPT-6 Astra (OpenAI, ChatGPT Pro mode).
\scanpage{119/8}
\mspage{6}

lines with four marked points or for elliptic
curves (!)) has never been properly made explicit, for example the relation
between the genus 0 case and the geometry of the octahedron and of the
tetrahedron. A fortiori, the moduli multiplicities $M_{0,5}$ (for
projective lines with five marked points) and $M_{1,2}$ (for
curves of genus 1 with two marked points), which are moreover almost
isomorphic to one another, seem virgin territory - braid groups
will not enlighten us about them ! I have begun to
look at $M_{0,5}$ in spare moments; it is a true jewel, with a
very rich geometry closely connected with that of the icosahedron.

The a priori interest of a complete knowledge of the first two
levels of the tower (namely, the cases in which the moduli dimension
$N = 3g-3+\nu$ is $\leq 2$) lies in this principle, that \uline{the entire tower
is reconstructed from the first two levels}, in the sense that, through
the fundamental ``gluing'' operation, level 1 supplies a
complete system of generators, and level 2 a complete system of
relations. There is a striking analogy, and I am convinced it is
more than formal, between this principle and Demazure's analogous
principle for the structure of reductive algebraic groups, when
the term ``level'' or ``moduli dimension'' is replaced by ``semisimple
rank of the reductive group''. The connection becomes still more striking
when one recalls that the Teichmüller group $T_{1,1}$ (now in the
discrete transcendental context, and not in the profinite algebraic
context, where one finds the profinite completions of the former) is
precisely $S\ell(2,\mathbb{Z})$, i.e. the group of integral points of the
``absolute'' simple group scheme of rank 1 $S\ell(2)_{\mathbb{Z}}$. Thus, \uline{the fundamental
building block for the Teichmüller tower} \uline{is essentially
the same as that for ``the tower'' of reductive groups of all ranks} -
a group of which one can probably say, moreover, that it is present
in all the essential disciplines of mathematics.

This principle for constructing the Teichmüller tower has
not been proved at present - but I have no doubt that it
is valid. It would follow (through a theory of dévissage of stratified
structures - in this instance the $M^{\wedge}_{g,\nu}$ - still to be written,
cf. par. 5) from an extremely plausible property of the open

\margin{8}

